\documentclass[
    language=english,
    doctype=article,
    institution=none,
]{../../omnilatex}

\RequirePackage{config/document-settings}
\addbibresource{../../bib/bibliography.bib}

\begin{document}

\maketitle

\section{Background}

Meridian Health Systems is a regional healthcare network operating 12
hospitals and 45 outpatient clinics across the Midwest United States.
Serving over 2.5 million patients annually, Meridian relies on its
technology infrastructure to support electronic health records (EHR),
medical imaging, telehealth services, and operational analytics.

By early 2025, Meridian's on-premises data centers---consisting of
approximately 2,400 physical servers and 18,000 virtual machines---had
reached critical capacity. Hardware refresh cycles were falling behind,
operating costs were escalating, and the organization faced increasing
difficulty meeting regulatory requirements for data availability and
disaster recovery.

\section{Problem Statement}

Meridian's legacy infrastructure presented four interconnected challenges:

\begin{enumerate}
    \item \textbf{Scalability Constraints.} Peak demand during flu season
        and public health emergencies overwhelmed existing capacity,
        resulting in degraded EHR performance and delayed access to
        patient records.
    \item \textbf{Rising Costs.} Annual infrastructure maintenance costs
        had grown to \$38 million, consuming 22\% of the IT budget and
        diverting resources from innovation initiatives.
    \item \textbf{Compliance Risk.} The organization's disaster recovery
        capabilities did not meet the 15-minute Recovery Time Objective
        (RTO) mandated by its Joint Commission accreditation.
    \item \textbf{Technical Debt.} Legacy systems running unsupported
        operating versions created security vulnerabilities and increased
        the risk of unplanned downtime.
\end{enumerate}

\section{Analysis}

\subsection{Current State Assessment}

A 90-day assessment conducted by the Enterprise Architecture team in Q1
2025 revealed:

\begin{itemize}
    \item Average server utilization was only 18\%, indicating massive
        over-provisioning.
    \item Storage growth was projected at 40\% annually, driven by
        medical imaging and genomics data.
    \item The existing disaster recovery site had a Recovery Time Objective
        of 4 hours---far exceeding the 15-minute target.
    \item Annual unplanned downtime costs were estimated at \$2.1 million
        per hour of outage.
\end{itemize}

\subsection{Options Evaluation}

Three migration strategies were evaluated:

\begin{table}[h]
    \centering
    \begin{tabular}{lccc}
        \toprule
        \textbf{Criterion} & \textbf{Lift and Shift} & \textbf{Replatform} & \textbf{Refactor} \\
        \midrule
        Cost (3-year)        & \$28M  & \$35M  & \$52M \\
        Timeline             & 12 mo  & 18 mo  & 30 mo \\
        Risk level           & Low    & Medium & High \\
        Scalability gain     & 2x     & 5x     & 10x+ \\
        Compliance alignment & Partial & Full   & Full \\
        Operational overhead & High   & Medium & Low \\
        \bottomrule
    \end{tabular}
    \caption{Comparison of migration strategies.}
\end{table}

The replatform strategy was selected as the optimal balance of cost,
risk, and long-term benefit. This approach would migrate workloads to
managed cloud services while preserving existing application architectures
where feasible.

\section{Solution}

\subsection{Architecture Design}

The target architecture was designed around the following principles:

\begin{itemize}
    \item \textbf{Hybrid Cloud.} Sensitive workloads (EHR, PHI) would
        remain in a private cloud environment; non-sensitive workloads
        (analytics, internal tools) would migrate to a public cloud
        provider.
    \item \textbf{Infrastructure as Code.} All infrastructure would be
        provisioned and managed through Terraform modules, enabling
        reproducibility and version control.
    \item \textbf{Zero Trust Security.} Network segmentation, identity-
        based access controls, and continuous monitoring would replace
        the traditional perimeter-based model.
    \item \textbf{Disaster Recovery.} Multi-region active-passive
        configuration with automated failover to achieve a 10-minute RTO.
\end{itemize}

\subsection{Migration Approach}

The migration was executed in four phases:

\begin{enumerate}
    \item \textbf{Foundation (Months 1--3).} Established cloud accounts,
        identity federation, network connectivity, and security
        baselines. Deployed monitoring and logging infrastructure.
    \item \textbf{Non-Critical Workloads (Months 3--6).} Migrated
        development environments, analytics platforms, and internal
        business applications. Validated performance and cost models.
    \item \textbf{Clinical Systems (Months 6--9).} Migrated EHR
        integrations, medical imaging archives, and telehealth platforms
        with dedicated cutover windows and rollback procedures.
    \item \textbf{Optimization (Months 9--12).} Decommissioned legacy
        data centers, optimized cloud resource utilization, and
        implemented auto-scaling policies.
\end{enumerate}

\section{Results}

\subsection{Performance Outcomes}

Post-migration metrics demonstrated significant improvements:

\begin{itemize}
    \item \textbf{EHR Response Time:} Reduced from 4.2 seconds to 0.8
        seconds, a 81\% improvement.
    \item \textbf{System Availability:} Increased from 99.2\% to 99.97\%,
        equivalent to 25 hours less downtime annually.
    \item \textbf{Recovery Time Objective:} Achieved 8-minute RTO, well
        below the 15-minute target.
    \item \textbf{Peak Capacity:} Auto-scaling handled a 300\% demand
        surge during the January 2026 flu season without performance
        degradation.
\end{itemize}

\subsection{Financial Outcomes}

\begin{table}[h]
    \centering
    \begin{tabular}{lrr}
        \toprule
        \textbf{Metric} & \textbf{Before} & \textbf{After} \\
        \midrule
        Annual infrastructure cost  & \$38M  & \$26M \\
        IT budget share (infra)      & 22\%   & 14\% \\
        Capital expenditure (annual) & \$12M  & \$ 3M \\
        Disaster recovery cost       & \$ 4M  & \$ 1.5M \\
        \bottomrule
    \end{tabular}
    \caption{Financial impact comparison.}
\end{table}

\subsection{Operational Outcomes}

\begin{itemize}
    \item Infrastructure provisioning time reduced from 6 weeks to 45
        minutes for standard environments.
    \item Security incident response time improved from 4 hours to 12
        minutes through automated detection and alerting.
    \item IT staff redeployed 15 engineers from maintenance tasks to
        innovation projects, including predictive analytics and patient
        engagement tools.
\end{itemize}

\section{Lessons Learned}

The following lessons emerged from the migration project and should inform
future infrastructure modernization efforts:

\begin{enumerate}
    \item \textbf{Invest Early in Foundational Capabilities.} The 3-month
        foundation phase---spent on identity, networking, and security
        ---saved significant rework during later migration phases.
    \item \textbf{Clinical Workloads Require Special Handling.} EHR and
        imaging systems demanded extensive performance testing and
        clinical staff involvement during cutover. Generic cloud migration
        playbooks were insufficient.
    \item \textbf{Cost Optimization Is Continuous.} Initial cloud costs
        were 15\% higher than projected until rightsizing and reserved
        instances were applied. Organizations must plan for ongoing cost
        management as a permanent operational discipline.
    \item \textbf{Change Management Is as Critical as Technology.} Success
        depended on clinician adoption. Dedicated training sessions and
        feedback loops with frontline staff were essential to smooth
        transition.
    \item \textbf{Hybrid Architectures Introduce Complexity.} Managing
        two environments increased operational overhead. Investing in
        unified observability tooling from the outset would have reduced
        this burden.
\end{enumerate}

\printbibliography[heading=none]

\end{document}
